\chapter{Conclusão}
\label{chp:conclusao}

Este trabalho analisou os benefícios e os impactos da unificação de sistemas corporativos a partir do caso do aplicativo InLight, desenvolvido para o BackOffice de compra de energia do Ambiente de Contratação Regulada da Light. O ponto de partida foi um problema comum em organizações que acumularam sistemas ao longo do tempo: informações distribuídas entre várias fontes, que exigem consulta e consolidação manuais. No caso estudado, essas fontes incluem os sistemas da CCEE, a plataforma Thunders, um sistema de lançamento de notas parcialmente integrado e controles em planilhas.

Para analisar esse problema, o Capítulo~\ref{sec:fundamentacao} reuniu os conceitos de Sistemas de Informação, sistemas legados, integração e unificação, gestão de processos, aceitação de tecnologia e o funcionamento da comercialização de energia no ambiente regulado. O Capítulo~\ref{sec:revisao} examinou estudos empíricos sobre os impactos, a implantação e a avaliação de sistemas integrados. O Capítulo~\ref{sec:proposta} aplicou esses referenciais ao caso: diagnosticou o cenário atual, descreveu o InLight, propôs um plano de implantação e definiu como seus resultados devem ser avaliados.

\section{Contribuições}

A contribuição prática do trabalho para a organização consiste no diagnóstico do cenário atual do BackOffice ACR, no plano de implantação do InLight e nos instrumentos de avaliação, que podem ser reaplicados em novos ciclos de melhoria do aplicativo. A contribuição acadêmica consiste em aplicar, de forma integrada, referenciais consolidados de avaliação de processos, de sucesso de SI, de benefícios de sistemas corporativos e de aceitação de tecnologia a um contexto pouco explorado pela literatura revisada: a unificação de informações em rotinas de \textit{back-office} de compra de energia no ambiente regulado brasileiro (Seção~\ref{sec:rev-posicionamento}).

\section{Limitações}

As principais limitações, detalhadas na Seção~\ref{sec:limitacoes}, são o estudo de um caso único, que permite generalização analítica, mas não estatística \citep{yin2015}; a avaliação realizada durante a fase de implantação, quando parte dos benefícios pode ainda não ter se manifestado; o número restrito de participantes; e a impossibilidade de analisar dados confidenciais, como os valores financeiros dos contratos.

\section{Trabalhos Futuros}

A partir dos resultados e das limitações deste estudo, sugerem-se os seguintes trabalhos futuros:

\begin{itemize}
    \item \textbf{Extensão ao mercado livre:} adaptar o InLight às rotinas do Ambiente de Contratação Livre, cujos contratos são negociados bilateralmente \citep{brasil2004decreto};
    \item \textbf{Integração automatizada com as fontes:} substituir o registro manual por integrações automatizadas com as fontes de dados, avaliando o uso de modelos de dados padronizados do setor ou a automação robótica de processos;
    \item \textbf{Avaliação longitudinal:} aplicar a medição dos indicadores e o questionáriono inicio da implementação e após um período maior de uso, para verificar se os benefícios se mantêm ou se ampliam \citep{markus2000learning};
    \item \textbf{Replicação:} aplicar os instrumentos de avaliação propostos a outras áreas da empresa ou a outras distribuidoras, permitindo a comparação entre casos \citep{yin2015}.
\end{itemize}

\section{Considerações Finais}

A unificação de sistemas não é apenas uma decisão tecnológica. A literatura e o caso analisado indicam que seus benefícios dependem de três condições: que a solução responda a problemas concretos do processo, que o processo seja revisto junto com a tecnologia e que os usuários percebam valor no novo modo de trabalho. O InLight foi concebido para atender à primeira condição; o plano de implantação e os instrumentos de avaliação propostos neste trabalho buscam assegurar as outras duas e permitir que seus resultados sejam medidos de forma objetiva.
